\setlength{\parindent}{0pt}
% ------------------------------------------------------------
\chapter{Design of XXX}
\label{chap:design}
% ------------------------------------------------------------
\section{Pipeline Overview}
\label{sec:pipeline-overview}
filler
\section{Expression Representation}
\label{sec:expr-representation}
filler
  \subsection{A Hash-Consed \NoCaseChange{$n$}-ary Circuit DAG}
  \label{sec:hash-consed-dag}
filler
  \subsection{Canonicalization: Flattening and Cancellation}
  \label{sec:canonicalization}
filler
  \subsection{Multiplicative Factoring}
  \label{sec:factoring}
filler
    \subsubsection{The Factoring Transformation}
    \label{sec:factoring-transformation}
filler
    \subsubsection{Enlarging the Class of Verifiable Circuits}
filler
    \subsubsection{Cost Analysis and the Factoring--Speed Tension}
    \label{sec:factoring-cost}
filler
\section{Component I: The Rewrite Engine}
\label{sec:rewrite-engine}

Given a chosen probe $\mathbf{w} = (w_1, \dots, w_k)$, $k \le d$, represented as node
IDs into the hash-consed circuit DAG of Section~\ref{sec:hash-consed-dag}, the rewrite
engine decides whether $\mathbf{w}$ is secret-independent. It does so constructively:
rather than reasoning about probability distributions directly, it searches for a
finite sequence of applications of Lemma~\ref{lem:substitution}
(Section~\ref{sec:randomness-substitution}) after which $\mathbf{w}$'s image is
\emph{syntactically} free of secret leaves, i.e.\ no secret leaf is reachable in the
sub-DAG below any $w_i$. Termination with a secret-free image certifies $\mathbf{w}$;
Theorem~\ref{thm:rewrite-soundness} in Section~\ref{sec:coupling-view} makes precise
why this certificate is sound.

The engine has three tiers of rewriting rules increasing in generality and cost,
summarized in
Algorithm~\ref{alg:rewrite-engine} and mirrored in the structure of this section.
\textproc{SimpleRewrite} (Section~\ref{sec:simple-rule}) handles the common case,
every random occurring at most once, with a cheap, purely reference-counted
procedure that never builds a new DAG node. It is tried first on the unfactored
probe and, on failure, once more after the multiplicative factoring of
Section~\ref{sec:factoring} has had a chance to expose additional single-occurrence
randoms. \textproc{GeneralRewrite} (Section~\ref{sec:general-substitution}) is the
complete fallback for the remaining case, a random occurring more than once: it
substitutes a random by its context wherever the random occurs, interleaving
factoring lazily whenever both of its own find-rules stall. Section~\ref{sec:coupling-view}
closes the section by reading the whole derivation as an explicit, replayable
bijection on the randomness assignment, the \gls{coupling} that Component III
(Section~\ref{sec:extension}) later replays against every other wire in the circuit.

\begin{algorithm}
\caption{Discharging a probe (Component I)}
\label{alg:rewrite-engine}
\begin{algorithmic}[1]
\Function{CheckProbe}{$\mathbf{w}$}
  \Require probe $\mathbf{w}$, node IDs into the circuit DAG $D$ (Section~\ref{sec:hash-consed-dag})
  \Ensure \textsc{Secure}, or \textsc{Insecure}$(\mathbf{w})$
  \State $(\mathit{ok}, T) \gets \Call{SimpleRewrite}{\mathbf{w}}$ \Comment{Tier 1: unfactored}
  \If{$\mathit{ok}$} \Return \textsc{Secure} \EndIf
  \State $(\mathit{ok}, T) \gets \Call{SimpleRewrite}{\mathrm{Factor}(\mathbf{w})}$ \Comment{Tier 2: factor (\S\ref{sec:factoring}), retry}
  \If{$\mathit{ok}$} \Return \textsc{Secure} \EndIf
  \State $(\mathit{ok}, T) \gets \Call{GeneralRewrite}{\mathbf{w}, \emptyset, \mathit{fuel}, \mathit{false}, \langle\rangle}$ \Comment{Tier 3: complete fallback}
  \State \Return $\mathit{ok}$ ? \textsc{Secure} : \textsc{Insecure}($\mathbf{w}$)
\EndFunction
\Statex
\Function{SimpleRewrite}{$\mathit{roots}$}
  \State $\mathit{todo} \gets \{\, r \in \mathrm{vars}(\mathit{roots}) : r \text{ occurs once, additively}, r \notin \mathit{roots} \,\}$, shallowest first
  \State $T \gets \langle\rangle$
  \While{$\mathit{todo} \neq \langle\rangle$ \textbf{and not} \Call{SecretFree}{$\mathit{roots}$}}
    \State $r \gets$ pop shallowest element of $\mathit{todo}$; \ $x \gets$ unique additive parent of $r$ \Comment{$x = e \oplus r$}
    \State mark $x$ as a fresh random leaf; \ $T \gets T.\mathrm{push}(r, x)$ \Comment{$t = e \oplus r = x$, already in hand}
    \State drop $r$; decrement occurrence counts of $x$'s other children
    \State \emph{(a child newly matching the pattern above, and $\notin \mathit{roots}$, joins $\mathit{todo}$)}
  \EndWhile
  \State \Return $(\Call{SecretFree}{\mathit{roots}}, T)$ \Comment{same shears, now recorded --- \S\ref{sec:coupling-view}}
\EndFunction
\Statex
\Function{GeneralRewrite}{$\mathit{tuple}, \mathit{used}, \mathit{fuel}, \mathit{factored}, T$}
  \If{\Call{SecretFree}{$\mathit{tuple}$}} \Return $(\mathit{true}, T)$ \EndIf
  \If{$\mathit{fuel} = 0$} \Return $(\mathit{false}, T)$ \EndIf
  \State $r \gets \Call{FindSimple}{\mathit{tuple}}$ \Comment{occurs once, additively, not a root}
  \If{$r = \bot$} \ $r \gets \Call{FindGeneral}{\mathit{tuple}, \mathit{used}}$ \Comment{occurs $\geq 2\times$; some additive parent free of $r$} \EndIf
  \If{$r = \bot$}
    \If{$\mathit{factored}$} \Return $(\mathit{false}, T)$ \EndIf
    \State \Return $\Call{GeneralRewrite}{\mathrm{Factor}(\mathit{tuple}), \mathit{used}, \mathit{fuel}, \mathit{true}, T}$ \Comment{\S\ref{sec:factoring-cost}, lazy retry}
  \EndIf
  \State $x \gets$ chosen additive parent of $r$; \ $e \gets \bigoplus(\mathrm{children}(x) \setminus \{r\})$ \Comment{side cond.\ $r \notin \mathrm{vars}(e)$}
  \State $t \gets e \oplus r$; \ $\mathit{tuple}' \gets \mathit{tuple}[r \mapsto t]$ \Comment{substitute everywhere --- Lemma~\ref{lem:substitution}}
  \State \Return $\Call{GeneralRewrite}{\mathit{tuple}', \mathit{used} \cup \{r\}, \mathit{fuel}-1, \mathit{false}, T.\mathrm{push}(r,t)}$
\EndFunction
\end{algorithmic}
\end{algorithm}

  \subsection{The Simple Rewrite Rule}
  \label{sec:simple-rule}

  The cheapest tier discharges the pattern Corollary~\ref{cor:single-occurrence-discharge}
  addresses directly: a random that occurs exactly once in the sub-DAG below the
  probe, as a child of an XOR node. A single depth-first traversal from the probe
  roots computes, for every node $n$ reachable from them, four pieces of bookkeeping:
  $\mathrm{totalParCount}(n)$, the number of edges into $n$ from within the explored
  region; $\mathrm{xorParCount}(n) \le \mathrm{totalParCount}(n)$, the number of those
  edges whose source is an XOR node; $\mathrm{mulDepth}(n)$, the number of AND-levels
  separating $n$ from the nearest root; and $\mathrm{parents}(n)$, the reverse
  adjacency needed to walk back up from a random to its context. A random $r$ is a
  \emph{simple-rule candidate} when $\mathrm{totalParCount}(r) =
  \mathrm{xorParCount}(r) = 1$ and $r$ is not itself a probe root: this is
  Corollary~\ref{cor:single-occurrence-discharge}'s hypothesis, checked without ever
  constructing $e$. Write $x$ for $r$'s unique parent, so $x = e \oplus r$ for the
  (unconstructed) context $e$.

  A probe root is excluded from candidacy because its value is what the adversary
  observes: it is fixed, not fresh, and relabeling it would claim independence for
  something the tuple actually exposes. The tuple
  $(r,\ s \oplus r,\ r \oplus r_3)$ makes the danger concrete. Its first entry is the
  bare leaf $r$, and the first two entries XOR to the secret $s$, so this tuple
  leaks $s$ and must be rejected. But once $r_3$ is discharged (it occurs once,
  additively, in $r \oplus r_3$), $r$'s own reference count drops to one: its only
  remaining parent is $s \oplus r$. Without the root guard, $r$ would now match the
  candidate pattern and be cascaded away, relabeling the very value the adversary
  observes and reporting a false \textsc{Secure}.

  Discharging a candidate (\textproc{ApplyRewrite} in the implementation) does not
  rebuild $x$ as $e \oplus r$; it marks $x$ itself as a fresh random leaf for every
  purpose the bookkeeping cares about, drops $r$, and decrements the reference counts
  of $x$'s other children, each of which loses the one incoming edge from $x$. This
  is licensed by Lemma~\ref{lem:substitution} applied with this specific $r$ and $e$:
  $x$, jointly with everything not depending on $r$, is distributed exactly as fresh,
  independent randomness, so nothing is lost by never inspecting $e$'s content.
  Because the DAG is hash-consed (Section~\ref{sec:hash-consed-dag}), $x$ may have
  several parents; a single discharge blinds every one of them at once, since all
  read the identical shared node.

  Decrementing a child's count can itself drop it to a value matching the candidate
  pattern, in which case it joins the work-list, ordered by increasing
  $\mathrm{mulDepth}$ so that the shallowest, closest-to-a-root candidates fire
  first, matching the increasing-depth discharge order used by \texttt{maskVerif}
  \parencite{10.1007/978-3-030-29959-0_15}. This cascade is again guarded by root
  membership. The loop succeeds once every probe root fails to be a bare secret leaf
  and every secret leaf in the circuit has zero recorded incoming edges within the
  current bookkeeping, a check answerable in $O(1)$ from the maintained counts, with
  no fresh traversal. It fails once the work-list empties without reaching that
  state. A random is discharged at most once, since its bookkeeping is erased the
  moment it fires, so the loop performs at most $|\mathcal{R}|$ rewrite steps.

% refer to the soundness proof
  This tier is sound, since each step is a direct instance of
  Corollary~\ref{cor:single-occurrence-discharge}, but incomplete: it has nothing to
  offer a random occurring twice, such as a mask reused across two cross-terms of a
  masked multiplication. Section~\ref{sec:general-substitution}'s fallback handles
  that case, at the price of the incremental, allocation-free bookkeeping this tier
  relies on.

  \subsection{Fallback: General Substitution}
  \label{sec:general-substitution}

  \textproc{GeneralRewrite} handles randoms that occur more than once by performing
  the substitution of Lemma~\ref{lem:substitution} literally, rather than only its
  single-occurrence case (Corollary~\ref{cor:single-occurrence-discharge}). It picks
  a random $r$ and an additive occurrence $x_1 = e \oplus r$ with
  $r \notin \mathrm{vars}(e)$, then rewrites every root by substituting
  $r \mapsto e \oplus r$ throughout. Re-canonicalization
  (Section~\ref{sec:canonicalization}) collapses the chosen occurrence $x_1$ to the
  bare leaf $r$ by $\mathbb{F}_2$-cancellation, exactly as in
  Corollary~\ref{cor:single-occurrence-discharge}, while every other occurrence of
  $r$ absorbs $e$ instead of vanishing, cancelling further if it happens to share
  structure with $e$. The loop succeeds once no secret leaf remains reachable from
  the rewritten roots.

  Each round recomputes both find-rules on the tuple's current form: a fresh
  depth-first traversal from the current roots, repeated every round rather than
  maintained incrementally as in Section~\ref{sec:simple-rule}.
  \textproc{FindSimple} looks first for the same single-occurrence, non-root pattern
  as Section~\ref{sec:simple-rule}, since a general substitution can turn a
  multi-occurrence random into a single-occurrence one; it prefers the shallowest
  candidate. Failing that, \textproc{FindGeneral} looks for any random, not yet used
  by a previous general step, occurring at least twice with some additive parent
  $x_1$ whose other children do not themselves depend on it, the side condition
  $r \notin \mathrm{vars}(e)$ that Lemma~\ref{lem:substitution} requires. It prefers
  the deepest candidate, the opposite bias from \textproc{FindSimple}; the
  implementation records no rationale for this choice beyond the ordering itself. A
  context $e$ is always read off an XOR node's other children: anything else would
  mean $e$ contains $r$, making $r \mapsto e \oplus r$ fail to be a bijection, so the
  implementation fails loudly there instead of silently admitting an unsound
  relabeling. The same traversal answers whether the tuple is secret-free, since a
  secret survives exactly when it is itself a root or has a recorded parent in the
  explored sub-DAG, so no second traversal is needed for that check.

  When both find-rules stall on the tuple's current form, the engine factors
  (Section~\ref{sec:factoring}) once and retries, rather than declaring failure
  immediately. Factoring can move a random from a purely multiplicative position,
  where neither find-rule can see it, into an additive one, but on its own it
  eliminates no secret, so it must be interleaved with substitution rather than run
  once up front or once at the end. Factoring is symmetrically not tried before the
  simple rule either: Algorithm~\ref{alg:rewrite-engine} always tries
  \textproc{SimpleRewrite} unfactored first, because factoring can destroy a
  single-occurrence opportunity the unfactored form exposes, so trying the cheap,
  unfactored tier first is not just faster on the common case but strictly more
  complete.

  \textproc{FindGeneral} is needed precisely when a random occurs additively in two
  places that do not telescope into one another under the simple rule. Suppose $r$
  occurs as an operand of $x_1 = e_1 \oplus r$ in one root and of $x_2 = e_2 \oplus r$
  in another, with $r \notin \mathrm{vars}(e_1) \cup \mathrm{vars}(e_2)$.
  Substituting through $x_1$, i.e.\ $r \mapsto e_1 \oplus r$, turns $x_1$ into
  $e_1 \oplus (e_1 \oplus r) = r$ by cancellation, as in
  Corollary~\ref{cor:single-occurrence-discharge}, while the same substitution turns
  $x_2$ into $e_2 \oplus e_1 \oplus r$. If $e_1 \oplus e_2$ is itself secret-free,
  both roots are secret-free after this one step; otherwise the loop continues,
  substituting further or factoring to expose a new opportunity. This is the shape of
  the case a masked multiplication typically produces, where a single random masks
  two different cross-terms: neither term's random is disposable on its own, but the
  substitution above can still cancel the secret-dependent part they share. Whether a
  given instance is discharged this way, and in how many rounds, depends entirely on
  how much of $e_1 \oplus e_2$ the rest of the derivation can still cancel; the loop
  offers no guarantee beyond trying. Even granting unlimited fuel, it only composes
  single-coordinate shears $\widehat{\sigma}$, so it can only certify
  secret-independence witnessed by a chain of substitutions, possibly interleaved
  with factoring. Completeness is proved only in the linear/affine case
  \textcolor{orange}{(nerede proof?)}; whether a genuinely irreducible counterexample
  exists beyond it is not settled by this thesis.

  Termination is only argued informally, not proved. Each general substitution
  retires its random into $\mathit{used}$ and is never repeated for it, bounding the
  number of general steps by $|\mathcal{R}|$; between general steps, simple
  substitutions strictly reduce the candidate pool; factoring is retried at most once
  per stall before the engine gives up on the current form. Interleaving factoring
  with two mutually triggering find-rules resists a clean potential-function
  argument, so a generous fuel counter backstops the loop instead. Exhausting fuel is
  treated as failure, which is sound (Theorem~\ref{thm:rewrite-soundness} is never
  invoked on a fuel-out) but not necessarily tight: a fuel-out reports the probe as
  though it were genuinely insecure, without exhibiting an actual distinguishing
  observation.

  \subsection{Relation to the Coupling-Construction View}
  \label{sec:coupling-view}

  A discharge does more than answer \textsc{Secure}/\textsc{Insecure}. Read the
  right way, it hands back the object Chapter~\ref{chap:couplings} asks for: a
  coupling for the probe. We state that object declaratively first, then show how
  the engine produces one.

  Let the secrets range now, rather than staying fixed as in
  Section~\ref{sec:randomness-substitution}: write
  $\mathbf{w}(\rho; \mathbf{a})$ for the value of the probe under randomness
  assignment $\rho \in \{0,1\}^{\mathcal{R}}$ and secret assignment
  $\mathbf{a} \in \{0,1\}^{\mathcal{S}}$, the public inputs held fixed.

  \begin{definition}[Coupling of a probe]
  \label{def:probe-coupling}
  A \emph{coupling} of a probe $\mathbf{w}$ is a family
  $\{\,\Phi_{\mathbf{a},\mathbf{b}}\,\}_{\mathbf{a},\mathbf{b} \in \{0,1\}^{\mathcal{S}}}$
  of measure-preserving bijections of $\{0,1\}^{\mathcal{R}}$, one per ordered
  pair of secret assignments, satisfying
  \[
    \mathbf{w}(\rho; \mathbf{a}) \;=\; \mathbf{w}\big(\Phi_{\mathbf{a},\mathbf{b}}(\rho);\, \mathbf{b}\big)
    \qquad \text{for every } \rho \in \{0,1\}^{\mathcal{R}} .
  \]
  \end{definition}

  Definition~\ref{def:probe-coupling} says precisely where to send the randomness
  so that a run under $\mathbf{a}$ is reproduced, observation for observation, by
  a run under $\mathbf{b}$: each $\Phi_{\mathbf{a},\mathbf{b}}$ dictates the
  relabelling of the random tape that absorbs the change of secret. Since it is
  measure-preserving, pushing the uniform tape through it turns this pointwise
  identity into an equality of laws, so the distribution of
  $\mathbf{w}(\cdot;\mathbf{a})$ matches that of $\mathbf{w}(\cdot;\mathbf{b})$. A
  coupling existing for \emph{every} pair $(\mathbf{a},\mathbf{b})$ is thus the
  same statement as $\mathbf{w}$ being secret-independent --- which is what makes
  it the right certificate. The family carries the structure one expects of such
  a certificate: one may take $\Phi_{\mathbf{a},\mathbf{a}} = \mathrm{id}$,
  $\Phi_{\mathbf{b},\mathbf{a}} = \Phi_{\mathbf{a},\mathbf{b}}^{-1}$, and
  $\Phi_{\mathbf{b},\mathbf{c}} \circ \Phi_{\mathbf{a},\mathbf{b}} =
  \Phi_{\mathbf{a},\mathbf{c}}$, so a single anchor
  $\Phi_{\mathbf{a}_0,\,\cdot}$ already generates all of it. It is these
  properties we care about, not any one procedure that realises them.

  What the engine records is one such procedure. Discharging a probe amounts to
  finding a \emph{derivation}: a finite sequence of rewrites
  \[
    \sigma_1 = (r_1 \mapsto e_1 \oplus r_1), \ \dots, \ \sigma_K = (r_K \mapsto e_K \oplus r_K),
  \]
  each valid --- $r_i \notin \mathrm{vars}(e_i)$ --- against the tuple's state
  just before step $i$, after which no secret leaf remains reachable. A rewrite
  is a rewrite: whatever rule proposed it, and whether the random it retires
  occurred once or many times, each $\sigma_i$ is a substitution in the sense of
  Definition~\ref{def:substitution}, hence a single measure-preserving shear
  $\widehat{\sigma_i}$ (Lemma~\ref{lem:substitution}). The coupling is assembled
  from these shears and from nothing else, so the construction below never needs
  to know how any individual rewrite was found.

  \begin{theorem}[Soundness of discharge]
  \label{thm:rewrite-soundness}
  If \textup{\textproc{CheckProbe}}$(\mathbf{w})$
  (Algorithm~\ref{alg:rewrite-engine}) returns \textup{\textsc{Secure}}, then the
  joint distribution of $\mathbf{w}(\rho; \mathbf{a})$, for $\rho$ sampled
  uniformly at random, does not depend on the secret assignment $\mathbf{a}$.
  Moreover the recorded substitutions realise a coupling of $\mathbf{w}$ in the
  sense of Definition~\ref{def:probe-coupling}.
  \end{theorem}

  \begin{proof}[Proof sketch]
  Fix a secret assignment $\mathbf{a}$. Each recorded step contributes the tape
  shear $\widehat{\sigma_i}$, whose defining context $e_i$ may itself read the
  secrets, so write $\widehat{\sigma_i}^{\,\mathbf{a}}$ for its action under
  $\mathbf{a}$ and set
  \[
    \varphi_{\mathbf{a}} \;=\; \widehat{\sigma_1}^{\,\mathbf{a}} \circ
    \widehat{\sigma_2}^{\,\mathbf{a}} \circ \cdots \circ \widehat{\sigma_K}^{\,\mathbf{a}},
  \]
  a composition of measure-preserving bijections, hence itself one. Writing
  $\mathbf{w}_K = \mathbf{w}\sigma_1 \cdots \sigma_K$ for the final tuple and
  chaining Lemma~\ref{lem:substitution} through the $K$ steps gives
  $\mathbf{w}\big(\varphi_{\mathbf{a}}(\rho); \mathbf{a}\big) = \mathbf{w}_K(\rho)$
  for every $\rho$. A \textsc{Secure} verdict means $\mathbf{w}_K$ is
  syntactically secret-free --- no secret leaf is reachable from it --- so
  $\mathbf{w}_K(\rho)$ does not depend on $\mathbf{a}$. The right-hand side is
  therefore common to all secret assignments, and for any two, $\mathbf{a}$ and
  $\mathbf{b}$,
  \[
    \mathbf{w}\big(\varphi_{\mathbf{a}}(\rho); \mathbf{a}\big) \;=\; \mathbf{w}_K(\rho)
    \;=\; \mathbf{w}\big(\varphi_{\mathbf{b}}(\rho); \mathbf{b}\big) .
  \]
  Substituting $\rho \mapsto \varphi_{\mathbf{a}}^{-1}(\rho)$ and setting
  $\Phi_{\mathbf{a},\mathbf{b}} = \varphi_{\mathbf{b}} \circ
  \varphi_{\mathbf{a}}^{-1}$ gives
  $\mathbf{w}(\rho; \mathbf{a}) = \mathbf{w}\big(\Phi_{\mathbf{a},\mathbf{b}}(\rho); \mathbf{b}\big)$
  for all $\rho$; the family $\{\Phi_{\mathbf{a},\mathbf{b}}\}$ is the promised
  coupling, and measure-preservation of each member yields the distributional
  claim. The chaining uses only that each $\sigma_i$ is a substitution
  (Definition~\ref{def:substitution}); a single-occurrence rewrite, where $e_i$
  is read off one $\oplus$-parent as in
  Corollary~\ref{cor:single-occurrence-discharge}, is folded in exactly like any
  other, so the argument is indifferent to how the derivation was produced.
  \end{proof}

  The composite $\varphi_{\mathbf{a}}$ is displayed above as a product of shears
  only for exhibition; it is never used. Component III (Section~\ref{sec:coupling-extension})
  uses a single member of the family as a \emph{blinding certificate}: it takes any
  other wire $u$ in the circuit and admits it into the certified \emph{safe set}
  when $u \circ \varphi_{\mathbf{a}}$ is secret-free --- when the same relabelling
  that blinds the probe blinds $u$ too.
  Evaluating $\varphi_{\mathbf{a}}$ on a sampled tape is a back-substitution
  through the recorded steps in reverse order $i = K, \dots, 1$, since the net
  image of $r_i$ is $t_i = e_i \oplus r_i$ evaluated against the assignment
  already updated by every later step: the derivation performs forward
  elimination, and reading the net map off is back-substitution through the same
  steps in reverse.

  The recorded target of a rewrite costs nothing to keep. The substitution
  $\sigma_i = (r_i \mapsto e_i \oplus r_i)$ sends $r_i$ to $t_i = e_i \oplus r_i$,
  and $t_i$ is exactly the $\oplus$-parent node the engine already held when it
  chose the rewrite --- there is no separate context $e_i$ to reconstruct. So
  every rewrite the engine makes, of whatever origin, can be appended to the
  derivation $\sigma_1, \dots, \sigma_K$ as it happens, and the coupling is read
  off that one list. Component III replays these shears directly against every
  other wire in the circuit; no probe and no rewrite is special-cased, and each
  step that helped certify the probe is a step that helps blind its neighbours.

\section{Component II: Traversing the Subset Space}
\label{sec:subset-space}
filler
  \subsection{The Space-Splitting Worklist Representation}
  \label{sec:space-splitting}
filler
  \subsection{Large-Set Pruning and Closure-Based Splitting}
  \label{sec:large-set-pruning}
filler
  \subsection{Completeness via Vandermonde's Identity}
  \label{sec:vandermonde}
filler
\section{Component III: Extending Probe Sets}
\label{sec:extension}
filler
  \subsection{Closures: Structural Deduction of Free Wires}
  \label{sec:closures}
filler
  \subsection{Coupling-Driven Extension}
  \label{sec:coupling-extension}
filler
